\documentclass[11pt,reqno,a4paper]{amsart}

\usepackage[T1]{fontenc}
\usepackage{lmodern}
\usepackage{amsmath,amssymb,mathtools}
\usepackage{booktabs}
\usepackage{array}
\usepackage{enumitem}
\usepackage{microtype}
\usepackage[hidelinks]{hyperref}

\newtheorem{theorem}{Theorem}[section]
\newtheorem{proposition}[theorem]{Proposition}
\newtheorem{lemma}[theorem]{Lemma}
\newtheorem{corollary}[theorem]{Corollary}
\theoremstyle{definition}
\newtheorem{definition}[theorem]{Definition}
\theoremstyle{remark}
\newtheorem{remark}[theorem]{Remark}

\newcommand{\comp}[1]{\overline{#1}}
\newcommand{\Cr}{\binom r2}
\newcolumntype{L}[1]{>{\raggedright\arraybackslash}p{#1}}

\title[The seven- and eight-color cases]{The seven- and eight-color cases of an Erd\H{o}s--Gy\'arf\'as balanced-coloring problem}
\author{Robert Sneiderman}
\date{20 July 2026}
\subjclass[2020]{05C15, 05C35, 05D10}
\keywords{edge coloring, balanced coloring, extremal graph theory, computer-assisted proof}

\begin{document}
\raggedbottom
\hypersetup{
  pdftitle={The seven- and eight-color cases of an Erdos--Gyarfas balanced-coloring problem},
  pdfauthor={Robert Sneiderman}
}

\begin{abstract}
We record computer-assisted proofs of two fixed cases of a problem of
Erd\H{o}s and Gy\'arf\'as. Every seven-coloring of the edges of $K_{50}$
has eight vertices whose induced edges omit a color, and every
eight-coloring of the edges of $K_{65}$ has nine vertices whose induced
edges omit a color. The seven-color proof uses a direct exhaustive lemma on
weighted graphs with seven vertices. The eight-color proof uses a human
full-color extremal argument together with 862 independently reconstructed
and LRAT-checked finite instances. The accompanying repository contains the
source, semantic verifiers, manifests, replay script, and proof artifacts.
These results are being circulated for independent checking and have not yet
received external mathematical review. Neither result settles the problem
for arbitrary $r$.
\end{abstract}

\maketitle

\section{Introduction}

Erd\H{o}s and Gy\'arf\'as asked whether the following holds for every
integer $r\geq3$: whenever the edges of $K_{r^2+1}$ are colored with $r$
colors, some $r+1$ vertices induce a complete graph on which at least one
color is absent \cite{ErdosGyarfas1999}. The question is also recorded as
Erd\H{o}s Problem 617 \cite{ErdosProblems617}.

This paper concerns the next two fixed parameters after the previously
known and separately treated cases.

\begin{theorem}\label{thm:combined}
Let $r\in\{7,8\}$. For every map
\[
  \chi:E(K_{r^2+1})\longrightarrow[r],
\]
there is a set $S\subseteq V(K_{r^2+1})$ with $|S|=r+1$ such that
\[
  \chi\!\left(E(K_{r^2+1}[S])\right)\ne[r].
\]
\end{theorem}

The two parts of Theorem~\ref{thm:combined} are proved by different finite
inputs. The $r=7$ finite lemma is a direct enumeration of
$12{,}618{,}770$ weighted cases, reconstructed a second way using
unlabelled graphs and automorphism orders. It does not use SAT. The $r=8$
proof reduces one branch to a single formula and another to 861
core--shell formulas. Independent semantic checkers reconstruct the
formulas, and an LRAT checker verifies their unsatisfiability.

\begin{remark}[Scope]\label{rem:scope}
Theorem~\ref{thm:combined} is a pair of fixed-parameter results. It does
not prove Erd\H{o}s Problem 617 for every $r$. It does not prove the fixed
case $r=9$, give a Lean formalization, or claim external review. The
certificate archives are not embedded in this manuscript.
\end{remark}

\section{Shared setup}\label{sec:setup}

Assume that every $(r+1)$-set sees all $r$ colors. For a fixed color $i$,
let $G_i$ be the graph formed by its edges. Then
\begin{equation}\label{eq:alpha}
  \alpha(G_i)\leq r.
\end{equation}
Every $(r+1)$-set contains at least one edge of each of the other $r-1$
colors, so
\begin{equation}\label{eq:local-cap}
  e(G_i[S])\leq\binom{r+1}{2}-(r-1)=\binom r2+1
  \qquad\left(S\in\binom{V(G_i)}{r+1}\right).
\end{equation}

For $n=ar+b$, where $0\leq b<r$, put
\begin{equation}\label{eq:pr}
  p_r(n)=r\binom a2+ab.
\end{equation}
This is the minimum number of edges in an $n$-vertex graph with
independence number at most $r$. Since the other $r-1$ color graphs
partition the complement of $G_i[W]$, Tur\'an's theorem gives the
full-color induced-density bound
\begin{equation}\label{eq:colored-density}
  e\!\left(\comp{G_i[W]}\right)\geq(r-1)p_r(|W|).
\end{equation}
This inequality uses coexistence of all colors; it is not a consequence of
the one-color local cap alone.

Choose a least color graph $G$. Its edge budget is
\begin{equation}\label{eq:least-budget}
  e(G)\leq M_r:=\frac{r(r^2+1)}2.
\end{equation}
The shared reduction gives
\begin{equation}\label{eq:degree-window}
  2\leq\delta(G)\leq r-1.
\end{equation}
The upper bound uses Brooks's theorem \cite{Brooks1941}; the lower bound
uses the Kang--Pikhurko extremal theorem \cite{KangPikhurko2005}.

\begin{proof}[Proof of \eqref{eq:degree-window}]
The average degree of $G$ is at most $r$. If $\delta(G)=r$, then $G$ is
$r$-regular. On the other hand,
\[
  \chi(G)\geq
  \left\lceil\frac{r^2+1}{\alpha(G)}\right\rceil\geq r+1.
\]
Some connected component therefore has chromatic number at least $r+1$.
Brooks's theorem makes that component $K_{r+1}$, contrary to
\eqref{eq:local-cap}.

For the lower bound, suppose that a vertex $v$ has degree
$d\in\{0,1\}$ and let $U$ be its nonneighbors other than $v$. Then
\[
  |U|=r^2-d,
  \qquad \alpha(G[U])\leq r-1.
\]
The complement of $G[U]$ is $K_r$-free. It is not $(r-1)$-partite:
such a partition would have a part of order at least $r+1$, which would
be a target-color $K_{r+1}$ in $G[U]$. The Kang--Pikhurko theorem, with
parameter $r-1$, gives
\[
 e(G[U])\geq
 p_{r-1}(|U|)+\left\lfloor\frac{|U|}{r-1}\right\rfloor-1.
\]
For $d=0$ the right side is $M_r+r+1$. For $d=1$ it is $M_r$ before
the edge incident with $v$ is added. Both cases contradict
\eqref{eq:least-budget}.
\end{proof}

We shall also use the strict colored-density form. If
$2\leq s<r$, $|W|=sr+t$, $t\geq1$, and
$\alpha(G_i[W])\leq s$, then
\begin{equation}\label{eq:nonpartite-density}
 e\!\left(\comp{G_i[W]}\right)
 \geq (r-1)\left(
 p_r(|W|)+\left\lfloor\frac{|W|}{r}\right\rfloor-1\right).
\end{equation}
Indeed, $G_i[W]$ is not $r$-colorable because $|W|>rs$. For every
$j\ne i$, it is a subgraph of $\comp{G_j[W]}$, so that complement is
not $r$-partite. It is $K_{r+1}$-free because $\alpha(G_j)\leq r$.
Applying the Kang--Pikhurko theorem to every $j\ne i$ and summing proves
\eqref{eq:nonpartite-density}.

\subsection{Colored terminal core and exact ladder}

\begin{theorem}[Colored terminal core]\label{thm:terminal-core}
Let $r\geq6$. There is no vertex set $W$ of order at least $2r$ such
that
\[
  \alpha(G_i[W])\leq2,
  \qquad \omega(G_i[W])\leq r-1.
\]
\end{theorem}

\begin{proof}
Put $H=G_i[W]$, $L=\comp H$, and $n=|W|$. The graph $L$ is
triangle-free and $\alpha(L)\leq r-1$. Every neighborhood in $L$ is
independent, whence
\[
  \Delta(L)\leq r-1,
  \qquad e(L)\leq\frac{n(r-1)}2.
\]
The colored-density inequality gives
\[
  e(L)\geq(r-1)p_r(n).
\]
At $n=2r$ the two bounds force $L$ to be $(r-1)$-regular. Since
$r-1>2(2r)/5$ for $r\geq6$, the Andr\'asfai--Erd\H{o}s--S\'os theorem
\cite{AndrasfaiErdosSos1974} makes $L$ bipartite. One side has order at least $r$, contrary to
$\alpha(L)\leq r-1$.

For $n>2r$, one has $p_r(n)>n/2$. This follows from $p_r(2r)=r$ and
\[
  p_r(m+1)-p_r(m)=\left\lfloor\frac mr\right\rfloor\geq2
  \qquad(m\geq2r).
\]
The two edge bounds then contradict one another.
\end{proof}

For fixed $r\geq6$, define $T_2=0$. Once $T_{s-1}$ is defined, put
\begin{align}
 A_s&=(s+1)r-2s-T_{s-1}+1,\label{eq:As}\\
 \widehat B_s&=rs(s+T_{s-1}-2)-2(r-1)(s-1),\label{eq:Bs}
\end{align}
and, when $A_s>0$ and the resulting integer is less than $r$, put
\begin{equation}\label{eq:Ts}
 T_s=\max\left\{1,
 \left\lfloor\frac{\widehat B_s}{A_s}\right\rfloor+1\right\}.
\end{equation}

\begin{theorem}[Exact colored core ladder]\label{thm:exact-ladder}
Whenever $T_s(r)$ is defined, there is no vertex set $W$ with
\[
 |W|\geq sr+T_s,
 \qquad \alpha(G_i[W])\leq s,
 \qquad \omega(G_i[W])\leq r-1.
\]
\end{theorem}

\begin{proof}
The case $s=2$ is Theorem~\ref{thm:terminal-core}. Assume the assertion
at $s-1$, and put $H=G_i[W]$ and $L=\comp H$. The graph $L$ is
$K_{s+1}$-free and has independence number at most $r-1$. For each
vertex $x$, its neighborhood in $L$ is $K_s$-free and retains that
independence bound. The induction hypothesis gives
\[
  \Delta(L)\leq(s-1)r+T_{s-1}-1=: \Delta_s.
\]
For $n=sr+t$, $1\leq t<r$, equation
\eqref{eq:nonpartite-density} gives
\[
  2e(L)\geq
  2(r-1)\left(r\binom s2+st+s-1\right).
\]
After subtracting the upper bound $n\Delta_s$, the result is
$A_st-\widehat B_s$. Definition \eqref{eq:Ts} is the first eligible
integer $t$ for which this is positive. For larger orders, the increment
of the lower bound minus that of $n\Delta_s$ is at least
\[
  2(r-1)\left\lfloor\frac nr\right\rfloor-\Delta_s
  \geq A_s>0.
\]
The contradiction therefore persists for every $n\geq sr+T_s$.
\end{proof}

In particular,
\begin{equation}\label{eq:threshold-values}
 T_3(r)=1\ (r\geq6),\qquad T_4(r)=2\ (r\geq6),
 \qquad (T_5(7),T_5(8))=(5,4).
\end{equation}
The same argument gives the pointwise form
\begin{equation}\label{eq:residual-min-degree}
 \delta(G_i[W])\geq r+t-T_{s-1}(r)
\end{equation}
whenever $|W|=sr+t$ and the independence and clique hypotheses at
level $s$ hold: the nonneighbors of a vertex have order at most
$(s-1)r+T_{s-1}-1$.

Let $v$ have color-$i$ degree $d\leq r-1$, and let $U$ be its nonneighbor
set. The colored core ladder proves the following block exclusion.

\begin{proposition}[Shared block exclusion]\label{prop:block-exclusion}
The graph $G_i[U]$ cannot contain $r-4$ pairwise disjoint copies of
$K_r$ in color $i$.
\end{proposition}

\begin{proof}
Since $v$ is anticomplete to $U$ in color $i$,
\begin{equation}\label{eq:U-alpha}
  \alpha(G_i[U])\leq r-1.
\end{equation}
Extend the proposed blocks to a maximal packing of $k$ target-color
copies of $K_r$ in $U$, let $R$ be the remainder, and put
$s=r-k-1$. Then
\begin{equation}\label{eq:R-alpha}
  \alpha(G_i[R])\leq s.
\end{equation}
For if $R$ contained an independent $(s+1)$-set, it could be extended
through the $k$ blocks by choosing one independent representative from
each. Every previously chosen vertex forbids at most one vertex in the
next $K_r$, by the local $(r+1)$-set cap, and
$(s+1)+k=r$. This would give an independent $r$-set in $U$, contrary to
\eqref{eq:U-alpha}. Maximality of the packing gives
\begin{equation}\label{eq:R-omega}
  \omega(G_i[R])\leq r-1.
\end{equation}

Write $t=r-d\geq1$. Since $|U|=r^2-d=r(r-1)+t$, the remainder has
order
\begin{equation}\label{eq:R-order}
  |R|=sr+t.
\end{equation}
If $k=r-4$, then $s=3$, and
\eqref{eq:R-alpha}--\eqref{eq:R-order} contradict
Theorem~\ref{thm:exact-ladder} because $T_3(r)=1$. If $k=r-3$, the
terminal-core theorem gives the same contradiction at $s=2$. If
$k=r-2$, then \eqref{eq:R-alpha} makes $R$ a clique, while
$|R|=r+t\geq r+1$, a contradiction. Finally, a packing with
$k\geq r-1$ contains $r-1$ blocks and leaves at least one vertex outside
them. That vertex and one independent representative from each block
give an independent $r$-set in $U$. These cases exhaust every maximal
extension of the proposed $r-4$ blocks.
\end{proof}

\subsection{Recursive colored edge floors}

For integers $a,m\geq0$, let $\mathcal F_r(a,m)$ be the family of
actual induced target-color graphs $G_i[W]$ with $|W|=m$ and
\[
  \alpha(G_i[W])\leq a,
  \qquad \omega(G_i[W])\leq r-1.
\]
Every member and each of its induced subgraphs satisfies
\[
  e(G_i[Z])\leq D_r(|Z|)
  :=\binom{|Z|}{2}-(r-1)p_r(|Z|).
\]
The full-color provenance in this definition is necessary; an abstract
one-color graph satisfying the local cap need not satisfy the terminal
exclusions above.

We define a lower bound $B_r(a,m)$, taking the value $+\infty$ when the
recursion proves that $\mathcal F_r(a,m)$ is empty. Put
\[
 B_r(0,0)=0,
 \qquad B_r(0,m)=+\infty\quad(m>0),
\]
and
\[
 B_r(1,m)=
 \begin{cases}
  \binom m2,&m\leq r-1,\\
  +\infty,&m\geq r.
 \end{cases}
\]
For $a\geq2$, define
\begin{equation}\label{eq:Qr}
 Q_r(a,m)=p_a(m)+
 \begin{cases}
  0,&m\leq a(r-1),\\
  \lfloor m/a\rfloor-1,&a(r-1)<m\leq ar,\\
  \lfloor m/a\rfloor,&m\geq ar+1,
 \end{cases}
\end{equation}
and
\begin{equation}\label{eq:Cr}
 C_r(\delta)=
 \begin{cases}
  \binom{\delta+1}{2},&0\leq\delta<r-1,\\
  \binom r2+1,&\delta=r-1,\\
  \displaystyle \delta+
  \left\lceil\frac{r+2}{r}\binom\delta2\right\rceil,&\delta\geq r.
 \end{cases}
\end{equation}
If $T_a(r)$ is defined and $m\geq ar+T_a(r)$, set
$B_r(a,m)=+\infty$. Otherwise put
\begin{equation}\label{eq:Rr}
 R_r(a,m)=
 \min_{\substack{0\leq\delta<m\\
 B_r(a-1,m-1-\delta)<+\infty}}
 \max\left\{
 B_r(a-1,m-1-\delta)+C_r(\delta),
 \left\lceil\frac{m\delta}{2}\right\rceil
 \right\},
\end{equation}
where the minimum of an empty set is $+\infty$. Put
\[
 b_r(a,m)=\max\{Q_r(a,m),R_r(a,m)\}.
\]
Set $B_r(a,m)=+\infty$ if $b_r(a,m)>D_r(m)$; otherwise set
$B_r(a,m)=b_r(a,m)$.

\begin{proposition}[Recursive colored lower bound]
\label{prop:recursive-lower-bound}
Every $F\in\mathcal F_r(a,m)$ satisfies
\[
  e(F)\geq B_r(a,m).
\]
In particular, $B_r(a,m)=+\infty$ certifies that the family is empty.
\end{proposition}

\begin{proof}
Induct on $a$. The cases $a=0,1$ are immediate. For $a\geq2$,
Tur\'an's theorem gives $e(F)\geq p_a(m)$. If $m>a(r-1)$, the
complement of $F$ is $K_{a+1}$-free and is not $a$-partite: an
$a$-partition would give a clique of order at least $r$ in $F$.
Kang--Pikhurko therefore gives the middle line of \eqref{eq:Qr}.

When $m\geq ar+1$, equality in that bound is incompatible with the local
cap. If $m\geq ar+2$, the old parts in the equality construction have
total order $m-1>ar$, so one part gives a clique of order at least
$r+1$. If $m=ar+1$, all old parts have order $r$. In the notation of
the equality construction, a special old part $S$, exceptional vertex
$x$, witness vertex $y$, and nonempty proper $A\subset S$ give
\[
 e(F[S\cup\{x\}])=\binom r2+r-|A|,
 \qquad
 e(F[S\cup\{y\}])=\binom r2+|A|.
\]
The local cap would require both $r-|A|\leq1$ and $|A|\leq1$, which
is impossible for $r\geq6$. This proves $e(F)\geq Q_r(a,m)$.

Choose a minimum-degree vertex $x$, set $\delta=d_F(x)$,
$N=N_F(x)$, and $U=V(F)\setminus(N\cup\{x\})$. Then
$F[U]\in\mathcal F_r(a-1,m-1-\delta)$. Put
\[
 X=e_F(N,U),\qquad Y=e_F(N).
\]
Minimum degree gives $X+2Y\geq\delta(\delta-1)$. As
$Y\leq\binom\delta2$, this first gives
$X+Y\geq\binom\delta2$. If $\delta=r-1$, equality would force
$Y=\binom{r-1}{2}$ and $X=0$, making $N\cup\{x\}$ a $K_r$;
one extra edge is therefore required.

If $\delta\geq r$, average the local cap over the sets
$\{x\}\cup S$ with $S\in\binom Nr$. Each $S$ spans at most
$\binom{r-1}{2}$ edges, and hence
\[
  Y\leq\frac{r-2}{r}\binom\delta2.
\]
Combining this with the degree-sum inequality gives
\[
  \delta+X+Y\geq C_r(\delta).
\]
The induction hypothesis on $F[U]$ proves the first term in the maximum
in \eqref{eq:Rr}; the degree sum gives
$e(F)\geq\lceil m\delta/2\rceil$. Taking both bounds at the same
$\delta$, then minimizing over possible minimum degrees, proves
\eqref{eq:Rr}. Theorem~\ref{thm:exact-ladder} justifies each terminal
$+\infty$ value, and $e(F)\leq D_r(m)$ justifies the final emptiness
test.
\end{proof}

For a least color graph $G_i$, let $v$ be a minimum-degree vertex of
degree $d\leq r-1$, and let $U$ be its nonneighbor set. After selecting
$j$ disjoint target copies of $K_r$ in $U$, let $R_j$ be their
complement. If $R_j$ has no further $K_r$, representative selection gives
\[
 R_j\in\mathcal F_r(r-1-j,r^2-d-jr).
\]
Indeed, an independent $(r-j)$-set in $R_j$ could be extended by one
representative from each deleted block. At every choice, fewer than $r$
previously selected vertices each forbid at most one block vertex, so a
representative remains. The resulting independent $r$-set in $U$
contradicts \eqref{eq:U-alpha}.

For the edge accounting, put $N=N_{G_i}(v)$,
$X=e_{G_i}(N,U)$, and $Y=e_{G_i}(N)$. Minimum degree gives
$X+2Y\geq d(d-1)$, while $Y\leq\binom d2$; hence
$X+Y\geq\binom d2$. Subtracting the edges incident with $v$, these
neighborhood edges, and the internal edges of the selected blocks gives
\begin{equation}\label{eq:entry-upper}
 e(R_j)\leq E_{r,d,j}:=
 M_r-d-\binom d2-j\binom r2.
\end{equation}
Consequently,
\begin{equation}\label{eq:entry-criterion}
 B_r(r-1-j,r^2-d-jr)>E_{r,d,j}
\end{equation}
forces every packing of $j$ blocks to extend by one block.

The exact margins $B_r-E_{r,d,j}$ needed below are
\begin{center}
\begin{tabular}{c@{\quad}c@{\quad}rrr}
\toprule
$r$&$d$&$j=0$&$j=1$&$j=2$\\
\midrule
7&0&57&$+\infty$&$+\infty$\\
7&1&42&$+\infty$&$+\infty$\\
7&2&25&$+\infty$&$+\infty$\\
7&3&18&28&$+\infty$\\
7&4&10&9&$+\infty$\\
7&5&5&4&$+\infty$\\
7&6&$-1$&$-2$&$-3$\\
\bottomrule
\end{tabular}
\qquad
\begin{tabular}{c@{\quad}rrrr}
\toprule
$d$&$j=0$&$j=1$&$j=2$&$j=3$\\
\midrule
0&60&104&$+\infty$&$+\infty$\\
1&44&72&$+\infty$&$+\infty$\\
2&35&41&$+\infty$&$+\infty$\\
3&24&23&$+\infty$&$+\infty$\\
4&17&16&$+\infty$&$+\infty$\\
5&8&7&6&$+\infty$\\
6&3&2&1&$+\infty$\\
7&$-4$&$-5$&$-6$&$-7$\\
\bottomrule
\end{tabular}
\end{center}
The left table is for $r=7$ and the right table for $r=8$. Every entry
needed for $d\leq5$ at $r=7$, and for $d\leq6$ at $r=8$, is positive
or infinite. Criterion \eqref{eq:entry-criterion} therefore forces,
respectively, three or four disjoint blocks, contrary to
Proposition~\ref{prop:block-exclusion}. Together with
\eqref{eq:degree-window}, this proves
\begin{equation}\label{eq:fixed-degree-reductions}
 \delta(G)=6\quad(r=7),
 \qquad \delta(G)=7\quad(r=8).
\end{equation}

\section{The seven-color case}\label{sec:r7}

\begin{theorem}\label{thm:r7}
Every seven-coloring of $K_{50}$ has an eight-vertex set that omits at
least one color.
\end{theorem}

The only new finite premise for this theorem is the next weighted
seven-vertex lemma.

\begin{lemma}[Weighted light edge]\label{lem:r7-light-edge}
Let $M$ be a graph on seven vertices with
\[
  \alpha(M)\leq5,
  \qquad 6\leq m:=e(M)\leq9.
\]
Give its vertices nonnegative integer weights $z_y$, and put
$Z=\sum_yz_y$.
\begin{enumerate}[label=\textup{(\roman*)}]
\item If $m+Z\leq8$, some edge $yy'$ satisfies
  \[
    d_M(y)+d_M(y')+z_y+z_{y'}\leq6.
  \]
\item If $m+Z\leq9$, some edge $yy'$ satisfies
  \[
    d_M(y)+d_M(y')+z_y+z_{y'}\leq7.
  \]
\end{enumerate}
\end{lemma}

The direct finite check gives the exact extrema
\begin{center}
\begin{tabular}{c@{\qquad}rrrr}
\toprule
$m\backslash Z$&0&1&2&3\\
\midrule
6&6&6&6&7\\
7&6&6&7&\\
8&6&6&&\\
9&7&&&\\
\bottomrule
\end{tabular}
\end{center}
over all cases satisfying the hypotheses. Section~\ref{sec:finite-r7}
states the semantics and independent reconstruction.

\subsection{Human implication chain}

Assume that every eight-set in a seven-coloring of $K_{50}$ sees all seven
colors, and choose a least color graph $G$.
\begin{enumerate}[label=\textup{R7.\arabic*},leftmargin=3.5em]
\item The least-color budget is $e(G)\leq175$.
\item The shared colored recursion excludes minimum degree at most five,
  while \eqref{eq:degree-window} gives the upper bound six. Hence
  $\delta(G)=6$.
\item For a degree-six vertex, its 43 nonneighbors $U$ satisfy
  \[
    \alpha(G[U])\leq6,
    \qquad e(G[U])\leq154.
  \]
\item Lemma~\ref{lem:r7-light-edge}, the star-and-cover argument, and the
  shared recursion give
  \[
    P_6(43)\geq156,
    \qquad P_5(36)\geq134,
    \qquad P_4(29)\geq113.
  \]
\item These floors exceed the successive edge upper bounds
  \[
    154,\qquad133,\qquad112.
  \]
  Three disjoint target copies of $K_7$ are therefore forced in $U$.
\item Proposition~\ref{prop:block-exclusion} forbids exactly three such
  blocks when $r=7$, giving the contradiction.
\end{enumerate}

The residual ledger is
\begin{center}
\begin{tabular}{c@{\qquad}c@{\qquad}c@{\qquad}c}
\toprule
residual order&independence cap&edge upper&no-next-$K_7$ floor\\
\midrule
43&6&154&156\\
36&5&133&134\\
29&4&112&113\\
\bottomrule
\end{tabular}
\end{center}

For this section, $P_a(n)$ denotes a valid lower edge bound on an actual
induced target-color graph $F$ with $|F|=n$, $\alpha(F)\leq a$, and
$\omega(F)\leq6$. The recursive bound above gives
\begin{equation}\label{eq:r7-input-floors}
\begin{gathered}
 P_3(20)\geq63,\qquad P_3(21)\geq75,\\
 P_5(37)\geq143,\qquad P_5(38)\geq153,
 \qquad P_5(39)\geq176.
\end{gathered}
\end{equation}

\subsection{The star-and-cover mechanism}

Let $x$ be a minimum-degree vertex of an induced target graph $F$, and
put
\[
 \delta=d_F(x),\qquad N=N_F(x),
 \qquad W=V(F)\setminus N_F[x].
\]
Write $M=\comp{F[N]}$ and $m=e(M)$. For $y\in N$, minimum degree
gives $d_F(y,W)\geq d_M(y)$. Define nonnegative integers $z_y$ by
\[
 d_F(y,W)=d_M(y)+z_y,
 \qquad Z=\sum_{y\in N}z_y.
\]
Then
\begin{equation}\label{eq:r7-cross-count}
 e_F(N,W)=2m+Z.
\end{equation}
If $\delta=6$, the number of edges outside $F[W]$ is
\begin{equation}\label{eq:r7-star6}
 6+e_F(N)+e_F(N,W)=21+m+Z.
\end{equation}
Here $M$ is nonempty, since otherwise $N_F[x]$ is a $K_7$. If
$\delta=7$, the corresponding total is
\begin{equation}\label{eq:r7-star7}
 28+m+Z.
\end{equation}
The local eight-set cap applied to $N_F[x]$ gives $m\geq6$. Also
$\alpha(M)\leq5$, since an independent six-set in $M$, together with
$x$, would be a target $K_7$.

For an edge $yy'$ of $M$, put
\[
 C=N_F(y,W)\cup N_F(y',W).
\]
Then
\begin{equation}\label{eq:r7-cover-size}
 |C|\leq d_M(y)+d_M(y')+z_y+z_{y'}.
\end{equation}
When $\alpha(F)\leq4$, the graph $F[W\setminus C]$ has independence
number at most two: an independent triple there, together with the
nonadjacent pair $y,y'$, would be an independent five-set in $F$.
Thus $|W\setminus C|\geq14$ contradicts
Theorem~\ref{thm:terminal-core}.

When $\delta=6$, choose any edge $yy'$ of the nonempty graph $M$. The
quantity $d_M(y)+d_M(y')-1$ counts the edges incident to at least one of
$y,y'$, so
\[
 d_M(y)+d_M(y')\leq m+1.
\]
Together with $z_y+z_{y'}\leq Z$, this bounds the right side of
\eqref{eq:r7-cover-size} by $m+1+Z$.

\begin{lemma}\label{lem:r7-p4-floors}
One has
\[
  P_4(28)\geq100,
  \qquad P_4(29)\geq113.
\]
\end{lemma}

\begin{proof}
Suppose first that $|F|=28$ and $e(F)\leq99$. Average degree gives
$\delta\leq7$. If $\delta\leq5$, then $|W|\geq22=3\cdot7+T_3(7)$,
contrary to Theorem~\ref{thm:exact-ladder}.

If $\delta=6$, then $|W|=21$ and
\eqref{eq:r7-input-floors} gives $e(F[W])\geq75$. Equation
\eqref{eq:r7-star6} gives $m+Z\leq3$; the bound following
\eqref{eq:r7-cover-size} gives $|C|\leq4$. Hence
$|W\setminus C|\geq17$, a contradiction.

If $\delta=7$, then $|W|=20$ and
$e(F[W])\geq63$. Equation \eqref{eq:r7-star7} gives $m+Z\leq8$.
Lemma~\ref{lem:r7-light-edge} gives $|C|\leq6$, so
$|W\setminus C|\geq14$, again a contradiction. This proves the first
floor.

Now suppose that $|F|=29$ and $e(F)\leq112$. Average degree gives
$\delta\leq7$. A degree at most six leaves at least 22 nonneighbors and
is excluded by the three-layer theorem. Thus $\delta=7$, $|W|=21$,
and $e(F[W])\geq75$. Equation \eqref{eq:r7-star7} gives
$m+Z\leq9$. Lemma~\ref{lem:r7-light-edge} gives $|C|\leq7$, leaving
at least 14 vertices and producing the terminal-core contradiction.
\end{proof}

\begin{lemma}\label{lem:r7-propagation}
One has
\[
 P_5(35)\geq122,
 \qquad P_5(36)\geq134,
 \qquad P_6(43)\geq156.
\]
\end{lemma}

\begin{proof}
If a 35-vertex graph with independence number at most five has at most
121 edges, its minimum degree is at most six. A degree at most four
leaves at least 30 nonneighbors and is excluded by $T_4(7)=2$. At
degrees five and six, Lemma~\ref{lem:r7-p4-floors} and the star floors
give
\[
 P_4(29)+15\geq128,
 \qquad P_4(28)+22\geq122.
\]
This proves $P_5(35)\geq122$.

For order 36 and at most 133 edges, the minimum degree is at most seven.
Degrees at most five are excluded by the four-layer theorem. At degrees
six and seven, the lower bounds are
\[
 P_4(29)+22\geq135,
 \qquad P_4(28)+34\geq134.
\]
This proves $P_5(36)\geq134$.

Finally, suppose that $|F|=43$, $\alpha(F)\leq6$, and
$e(F)\leq155$. The minimum degree is at most seven. A degree at most two
leaves at least 40 nonneighbors and is excluded by $T_5(7)=5$. For
degrees three through seven the lower bounds are
\[
\begin{array}{c|ccccc}
\delta&3&4&5&6&7\\ \hline
e(F)&176+6&153+10&143+15&134+22&122+34.
\end{array}
\]
Every entry is at least 156, proving the final floor.
\end{proof}

\begin{proof}[Proof of Theorem~\ref{thm:r7}]
Suppose otherwise and choose a least color graph $G$. Equation
\eqref{eq:fixed-degree-reductions} gives $\delta(G)=6$. Fix a
degree-six vertex $v$, and let $U$ be its 43 nonneighbors. Then
$\alpha(G[U])\leq6$. Minimum-degree accounting gives at least
$6+\binom62=21$ target edges outside $U$, and hence
\[
 e(G[U])\leq175-6-\binom62=154.
\]
If $G[U]$ had no $K_7$, Lemma~\ref{lem:r7-propagation} would give at
least 156 edges. Choose a target $K_7$, delete it, and call the
36-vertex remainder $R_1$. If $R_1$ had an independent six-set, each
of its vertices would forbid at most one vertex of the deleted $K_7$, so
one block vertex could be adjoined to form an independent seven-set in
$U$. Hence $\alpha(G[R_1])\leq5$, while
\[
 e(G[R_1])\leq154-\binom72=133.
\]
Lemma~\ref{lem:r7-propagation} forces a second disjoint $K_7$. Delete
it and call the 29-vertex remainder $R_2$. An independent five-set in
$R_2$ could be extended first through one deleted block and then through
the other, with at most five and then six forbidden block vertices. Thus
\[
 \alpha(G[R_2])\leq4,
 \qquad e(G[R_2])\leq133-\binom72=112.
\]
Lemma~\ref{lem:r7-p4-floors} forces a third disjoint $K_7$. Thus $U$
contains three disjoint target copies of $K_7$, contrary to
Proposition~\ref{prop:block-exclusion}.
\end{proof}

\section{Finite verification for the seven-color case}\label{sec:finite-r7}

A finite case is a pair $(M,z)$, where $M$ is a labeled graph on seven
vertices, $6\leq e(M)\leq9$, $\alpha(M)\leq5$, and
$z\in\mathbb N^7$ satisfies $e(M)+\sum_yz_y\leq9$. For each case, the
checker computes
\[
  \min_{yy'\in E(M)}
  \left(d_M(y)+d_M(y')+z_y+z_{y'}\right).
\]
The exact enumeration is as follows.

\begin{center}
\begin{tabular}{c@{\qquad}r@{\qquad}c}
\toprule
$(m,Z)$&weighted cases&largest possible minimum\\
\midrule
$(6,0)$&54,257&6\\
$(6,1)$&379,799&6\\
$(6,2)$&1,519,196&6\\
$(6,3)$&4,557,588&7\\
$(7,0)$&116,280&6\\
$(7,1)$&813,960&6\\
$(7,2)$&3,255,840&7\\
$(8,0)$&203,490&6\\
$(8,1)$&1,424,430&6\\
$(9,0)$&293,930&7\\
\midrule
total&12,618,770&\\
\bottomrule
\end{tabular}
\end{center}

The admissible labeled graph counts for $m=6,7,8,9$ are respectively
\[
  54{,}257,\quad116{,}280,\quad203{,}490,\quad293{,}930.
\]
The canonical checker enumerates labeled edge subsets and weak
compositions directly. The separate Sage checker uses nauty
\cite{McKayPiperno2014} to generate
40, 65, 97, and 131 unlabelled graph types, then reconstructs the labeled
counts from automorphism orders. Both routes reproduce the table. Neither
route calls a SAT solver, and there is no detached proof trace.

The accompanying release contains the direct checker
\path{code/verify_r7_weighted_light_edge.py}, the independent Sage checker
\path{code/verify_r7_weighted_light_edge_sage.py}, the arithmetic replay
\path{code/verify_r7_fixed_outer_closure.py}, and the shared ladder check
\path{code/verify_colored_core_ladder.py}. Their hashes are recorded in the
release checksum ledger.

\section{The eight-color case}\label{sec:r8}

\begin{theorem}\label{thm:r8}
Every eight-coloring of $K_{65}$ has a nine-vertex set that omits at least
one color.
\end{theorem}

Throughout this section, $P_a(n)$ denotes a valid lower edge bound on an
actual induced target-color graph with $n$ vertices, independence number
at most $a$, and clique number at most seven. The proof needs two such
edge floors.

\begin{lemma}[Twenty-three-vertex full-color floor]\label{lem:r8-p323}
Let $H$ be an actual induced target-color graph in a hypothetical balanced
eight-coloring. If
\[
  |V(H)|=23,
  \qquad\alpha(H)\leq3,
  \qquad\omega(H)\leq7,
\]
then $e(H)\geq91$.
\end{lemma}

The proof reduces a putative graph with at most 90 edges to a degree-seven
vertex with a 15-vertex triangle-free complement core. Full-color density
and the Kang--Pikhurko theorem force the extremal core. The missing-edge
graph on the seven-vertex neighborhood is a star, and each center--leaf
row pair partitions a minimum cover of the core. The full-color ten-set cap
removes one split core; the local nine-set cap removes the remaining two.

\begin{proof}
Suppose instead that $e(H)\leq90$, and choose a minimum-degree vertex
$v$. Average degree gives $d_H(v)\leq\lfloor180/23\rfloor=7$. If
$d_H(v)\leq6$, its nonneighbors contain at least 16 vertices. Their
induced graph has independence number at most two and clique number at
most seven, contrary to Theorem~\ref{thm:terminal-core}. Hence
$d_H(v)=7$.

Put $A=N_H(v)$ and $B=V(H)\setminus(A\cup\{v\})$, so $|A|=7$ and
$|B|=15$. Define
\[
 L=\comp{H[B]},\qquad M=\comp{H[A]},
 \qquad m=e(L),\qquad f=e(M).
\]
The graph $L$ is triangle-free, has $\alpha(L)\leq7$, and is
nonbipartite. Indeed, one side of a bipartition of 15 vertices would have
order at least eight and would give a target $K_8$ in $H[B]$.
The colored-density bound gives
\begin{equation}\label{eq:r8-p323-m-lower}
 m\geq7p_8(15)=49.
\end{equation}
The Kang--Pikhurko bound for a nonbipartite triangle-free graph on 15
vertices gives
\begin{equation}\label{eq:r8-p323-m-upper}
 m\leq t_2(15)-\left\lfloor\frac{15}{2}\right\rfloor+1
 =56-7+1=50.
\end{equation}
Here $t_2(15)=\lfloor15^2/4\rfloor=56$ is the triangle-free Tur\'an
number.

For $a\in A$, put
\[
 D_a=N_H(a)\cap B,
 \qquad c_a=|D_a|,
 \qquad c=\sum_{a\in A}c_a.
\]
Minimum degree gives
\[
 d_H(a)=1+6-d_M(a)+c_a\geq7,
 \qquad c_a\geq d_M(a),
 \qquad c\geq2f.
\]
Also $f\geq1$, since otherwise $\{v\}\cup A$ is a target $K_8$.
The full edge ledger is
\begin{equation}\label{eq:r8-p323-edge-ledger}
 e(H)=7+(21-f)+(105-m)+c=133-f-m+c.
\end{equation}
It follows that
\begin{equation}\label{eq:r8-p323-c-upper}
 c\leq m+f-43\leq f+7,
 \qquad f\leq7.
\end{equation}

For every edge $aa'\in E(M)$, the set $D_a\cup D_{a'}$ is a vertex
cover of $L$. Otherwise an uncovered edge of $L$, together with $a,a'$,
would be an independent four-set in $H$. Since
$\tau(L)=15-\alpha(L)\geq8$, we have
\begin{equation}\label{eq:r8-p323-cover}
 c_a+c_{a'}\geq8\qquad(aa'\in E(M)).
\end{equation}
The graph $M$ has no two disjoint edges, because they would give
$c\geq16$, while \eqref{eq:r8-p323-c-upper} gives $c\leq14$. If $M$
contained a triangle, the absence of two disjoint edges would make that
triangle its entire nonisolated part, so $f=3$. Summing
\eqref{eq:r8-p323-cover} around the triangle would give $c\geq12$, while
\eqref{eq:r8-p323-c-upper} would give $c\leq10$.

A nonempty graph with matching number at most one and no triangle is a
star. Thus
\[
 M=K_{1,f}\cup(6-f)K_1,
 \qquad1\leq f\leq6.
\]
Let $a_0$ be the center, $a_1,\ldots,a_f$ the leaves, and put
$x=c_{a_0}$. Minimum degree and \eqref{eq:r8-p323-cover} give
\[
 x\geq f,
 \qquad c_{a_j}\geq\max\{1,8-x\}\qquad(1\leq j\leq f).
\]
Consequently,
\begin{equation}\label{eq:r8-p323-c-lower}
 c\geq x+f\max\{1,8-x\}\geq f+7.
\end{equation}
For $x\leq7$, the difference between the final two quantities is
$(f-1)(7-x)$; for $x\geq7$, it is $x-7$. Equations
\eqref{eq:r8-p323-m-upper}, \eqref{eq:r8-p323-c-upper}, and
\eqref{eq:r8-p323-c-lower} force
\begin{equation}\label{eq:r8-p323-equality}
 m=50,\qquad c=f+7,\qquad e(H)=90.
\end{equation}

For every leaf $a_j$, equality also gives
\[
 c_{a_0}+c_{a_j}=8,
 \qquad D_{a_0}\cap D_{a_j}=\varnothing,
\]
and $C_j:=D_{a_0}\cup D_{a_j}$ is an eight-vertex minimum cover of
$L$. For $f=1$ this follows from $c=8$ and
\eqref{eq:r8-p323-cover}. For $f\geq2$, equality in
\eqref{eq:r8-p323-c-lower} forces $x=7$, every leaf row to have size
one, and every isolated row to be empty.

The equality case of the Kang--Pikhurko theorem now applies. There are
disjoint sets $P,Q$ and vertices $x_0,y_0$ such that
\[
 |P|=6,
 \quad |Q|=7,
 \quad Q=X\mathbin{\dot\cup}Y,
 \quad a:=|X|\in\{1,2,3\},
\]
and
\[
 E(L)=E(K_{P,Q})\cup\{x_0y_0\}
 \cup\{x_0q:q\in X\}\cup\{y_0q:q\in Y\}.
\]
Swapping $x_0,y_0$ permits $a\leq3$. The maximum independent
seven-sets and their complementary minimum covers are exactly
\begin{center}
\begin{tabular}{p{0.29\textwidth}p{0.38\textwidth}c}
\toprule
independent set&minimum cover $C$&$e_L(C)$\\
\midrule
$Q$&$P\cup\{x_0,y_0\}$&1\\
$P\cup\{x_0\}$&$Q\cup\{y_0\}$&$7-a$\\
$P\cup\{y_0\}$&$Q\cup\{x_0\}$&$a$\\
$\{x_0\}\cup Y$&$P\cup X\cup\{y_0\}$&6, if $a=1$\\
\bottomrule
\end{tabular}
\end{center}
This list is exhaustive. An independent set meeting $P$ cannot meet
$Q$, and order seven then forces all six vertices of $P$ and one of
$x_0,y_0$. An independent set avoiding $P$ is checked directly from
$Q=X\mathbin{\dot\cup}Y$.

Fix a leaf $a_j$. The ten-set
$S=\{a_0,a_j\}\cup C_j$ contains no target edge between $a_0,a_j$,
exactly eight cross edges, and $28-e_L(C_j)$ target edges inside $C_j$.
The full-color ten-set density bound, not merely the local nine-set cap,
therefore gives
\begin{equation}\label{eq:r8-ten-set}
 36-e_L(C_j)=e_H(S)\leq D_8(10)=31,
 \qquad e_L(C_j)\geq5.
\end{equation}

If $a=3$, every minimum cover in the table spans at most four edges,
contrary to \eqref{eq:r8-ten-set}. If $a=2$, the only remaining cover is
$C_j=Q\cup\{y_0\}$. Put $C'=Q\cup\{x_0\}$; then
$e_L(C')=2$ and $|C_j\cap C'|=7$. For each
$z\in\{a_0,a_j\}$, the local cap on the nine-set $\{z\}\cup C'$ gives
\[
 28-2+|D_z\cap C'|\leq29,
 \qquad |D_z\cap C'|\leq3.
\]
The two rows partition $C_j$, so their intersections with $C'$ sum to
seven, not at most six.

Finally, let $a=1$. If $C_j=Q\cup\{y_0\}$, use
$C'=Q\cup\{x_0\}$. If $C_j=P\cup X\cup\{y_0\}$, use
$C'=P\cup\{x_0,y_0\}$. In both cases
\[
 e_L(C')=1,
 \qquad |C_j\cap C'|=7.
\]
The local nine-set cap gives $|D_z\cap C'|\leq2$ for each of the two
rows, again contradicting their total intersection of seven. All equality
cores are impossible.
\end{proof}

The second floor is the computer-assisted dependency
\begin{equation}\label{eq:r8-p324}
  P_3(24)\geq106.
\end{equation}
Its exhaustive minimum-degree split is
\begin{center}
\small
\begin{tabular}{c p{0.36\textwidth}p{0.34\textwidth}}
\toprule
minimum degree&reduction&receipt\\
\midrule
$0,\ldots,4$&complement degree count&human arithmetic\\
5&one 18-vertex endpoint formula&semantic reconstruction and one LRAT\\
6&17-vertex complement degree argument&human proof\\
7&16-vertex complement--shell incidence argument&human proof\\
8&861 fixed core--shell pairs&semantic reconstruction and 861 LRATs\\
\bottomrule
\end{tabular}
\end{center}

\begin{proof}[Proof of \eqref{eq:r8-p324}]
Suppose that $H$ is a counterexample: $|V(H)|=24$,
$\alpha(H)\leq3$, $\omega(H)\leq7$, every nine-set spans at most 29
edges, and $e(H)\leq105$. Choose a minimum-degree vertex $v$ of degree
$d\leq8$, let $U$ be its $u=23-d$ nonneighbors, and put
$L=\comp{H[U]}$. The graph $L$ is triangle-free,
$\alpha(L)\leq7$, every nine-set of $L$ spans at least seven edges,
and $\Delta(L)\leq7$.

Minimum-degree accounting gives
\begin{equation}\label{eq:p83-basic-lower}
 e(H)\geq\binom{d+1}{2}+\binom u2-e(L).
\end{equation}
Using $e(L)\leq\lfloor7u/2\rfloor$ gives
\begin{center}
\begin{tabular}{c@{\quad}c@{\quad}c@{\quad}l}
\toprule
$d$&$u$&lower bound for $e(H)$&conclusion\\
\midrule
0&23&173&impossible\\
1&22&155&impossible\\
2&21&140&impossible\\
3&20&126&impossible\\
4&19&115&impossible\\
5&18&105&equality endpoint\\
\bottomrule
\end{tabular}
\end{center}

At $d=5$, equality in \eqref{eq:p83-basic-lower} forces
$e(L)=63$, so $L$ is seven-regular. Equality in the star accounting
forces $H[N_H[v]]$ to be a $K_6$ with no edges from $N_H(v)$ to $U$.
Thus the branch requires a triangle-free seven-regular graph on 18
vertices with independence number at most seven. The endpoint formula in
Section~\ref{sec:finite-r8} encodes exactly that graph and has a checked
LRAT refutation.

If $d=6$, then $u=17$. In fact $\Delta(L)\leq6$. To see this, suppose
that $z$ has seven neighbors $A$ in $L$; $A$ is independent. Among the
remaining nine vertices $b$, put $p(b)=|N_L(b)\cap A|$. For distinct
$b,c$, the nine-set $A\cup\{b,c\}$ gives
\[
 p(b)+p(c)+\mathbf1_{bc\in E(L)}\geq7.
\]
At most two vertices have $p(b)\leq3$: two such vertices must both have
value three and be adjacent, while three would form a triangle. Any two
vertices with $p$-value at least four are nonadjacent, because their
neighborhoods in the seven-set $A$ intersect. At least seven such
vertices, together with $z$, form an independent eight-set, a
contradiction. Hence $e(L)\leq51$, and
\[
 e(H)\geq\binom72+\binom{17}{2}-51=106.
\]

It remains to exclude $d=7,8$.

\medskip
\noindent\textit{The degree-seven branch.}
Here $|A|=7$ and $|B|=16$ for
$A=N_H(v)$ and $B=V(H)\setminus(A\cup\{v\})$. Put
$L=\comp{H[B]}$. It is triangle-free, has $\alpha(L)\leq7$, every
nine vertices span at least seven edges, and $\Delta(L)\leq7$.

We first prove $e(L)\leq48$. It is enough to show that a degree-seven
vertex forces $e(L)=44$. Suppose $z$ has degree seven, put
$P=N_L(z)$, and let $Q$ be the other eight vertices. The set $P$ is
independent. For $b\in Q$, write $p_b=|N_L(b)\cap P|$. For distinct
$b,c\in Q$,
\begin{equation}\label{eq:d7-pair}
 p_b+p_c+\mathbf1_{bc\in E(L)}\geq7.
\end{equation}
At most two vertices are low, meaning $p_b\leq3$. Any two low vertices
must have value three and be adjacent, while three low vertices would
form a triangle. Any two high vertices are nonadjacent, since their
neighborhoods in $P$ both have order at least four. There cannot be
seven high vertices, since those vertices together with $z$ would be an
independent eight-set. Thus $Q$ consists of six high vertices and two
low vertices $x,y$, with
\[
 p_x=p_y=3,
 \qquad xy\in E(L).
\]
The only possible edges inside $Q$ are $xy$ and edges from $x$ or $y$
to the six high vertices. No high vertex can meet both $x$ and $y$, so
$e(L[Q])\leq7$. The nine-set $\{z\}\cup Q$ gives the reverse
inequality, and every high vertex meets exactly one of $x,y$. Let $R$ be
the six high vertices, and put
$R_x=N_L(x)\cap R$ and $R_y=N_L(y)\cap R$. Then $R_x,R_y$ partition
$R$, and the degree bound gives $|R_x|=|R_y|=3$. For $r\in R_x$, the sets
$N_L(r)\cap P$ and $N_L(x)\cap P$ are disjoint; their orders are at
least four and exactly three, so they partition $P$. The same holds on
the $y$ side. Consequently
\[
 \sum_{b\in Q}p_b=2\cdot3+6\cdot4=30,
 \qquad e(L)=7+30+7=44.
\]
If no degree-seven vertex exists, $e(L)\leq16\cdot6/2=48$.

Put $F=\comp{H[A]}$, $f=e(F)$, and, for $a\in A$, put
\[
 D_a=N_H(a)\cap B,
 \qquad c_a=|D_a|,
 \qquad c=\sum_{a\in A}c_a.
\]
The graph $F$ is nonempty, or $H[\{v\}\cup A]$ is a $K_8$.
Minimum degree gives $c_a\geq d_F(a)$ and $c\geq2f$. The full edge
count is
\[
 e(H)=148-m-f+c\leq105,
 \qquad m=e(L),
\]
so
\begin{equation}\label{eq:d7-c-upper}
 c\leq m+f-43\leq f+5.
\end{equation}
For every $aa'\in E(F)$, the set $D_a\cup D_{a'}$ covers $L$; an
uncovered $L$-edge would make an independent four-set with $a,a'$.
Since $\tau(L)\geq9$,
\begin{equation}\label{eq:d7-cover}
 c_a+c_{a'}\geq9\qquad(aa'\in E(F)).
\end{equation}
Equations \eqref{eq:d7-c-upper}--\eqref{eq:d7-cover} give $f\geq4$,
while $c\geq2f$ and \eqref{eq:d7-c-upper} give $f\leq5$.

The graph $F$ cannot have two disjoint edges, since
\eqref{eq:d7-cover} would give $c\geq18$, whereas
\eqref{eq:d7-c-upper} gives $c\leq10$. A graph with matching number at
most one is a star or a triangle together with isolated vertices. Since
$f\in\{4,5\}$, it is a star $K_{1,f}$. If $a_0$ is its center and
$x=c_{a_0}$, then
\[
 c\geq x+f\max\{1,9-x\}\geq f+8,
\]
contrary to \eqref{eq:d7-c-upper}. This excludes $d=7$.

\medskip
\noindent\textit{The degree-eight branch.}
Now $|A|=8$ and $|B|=15$. Define
\[
 L=\comp{H[B]},\quad F=\comp{H[A]},
 \quad m=e(L),\quad f=e(F).
\]
The graph $L$ is triangle-free, $\alpha(L)\leq7$,
$\Delta(L)\leq7$, and every nine-set spans at least seven edges. The
graph $F$ is $K_4$-free and $\alpha(F)\leq6$. Applying the local cap
to $\{v\}\cup A$ gives $f\geq7$.

For $a\in A$, write
\[
 D_a=N_H(a)\cap B,
 \qquad c_a=d_F(a)+\varepsilon_a,
 \qquad \varepsilon_a\geq0.
\]
The equality follows from minimum degree. With
$c=\sum_ac_a=2f+\sum_a\varepsilon_a$, the edge ledger is
\[
 e(H)=8+(28-f)+(105-m)+c=141-f-m+c\leq105.
\]
Hence
\begin{equation}\label{eq:d8-slack}
 c\leq m+f-36,
 \qquad \sum_a\varepsilon_a\leq m-f-36.
\end{equation}
The first inequality and $f\geq7$ give $m\geq43$. The graph $L$ is
nonbipartite, because a bipartition has a part of order at least eight.
Kang--Pikhurko gives $m\leq50$, so
\begin{equation}\label{eq:d8-m-range}
 43\leq m\leq50.
\end{equation}

Two incidence conditions are forced. If $aa'\in E(F)$, the rows
$D_a,D_{a'}$ cover $L$, and therefore
\begin{equation}\label{eq:d8-pair-cover}
 d_F(a)+d_F(a')+\varepsilon_a+\varepsilon_{a'}\geq8.
\end{equation}
If $T$ is a triangle of $F$, the three corresponding rows cover all of
$B$, and therefore
\begin{equation}\label{eq:d8-triangle-cover}
 \sum_{a\in T}\bigl(d_F(a)+\varepsilon_a\bigr)\geq15.
\end{equation}

Nauty 2.8.9's \texttt{geng} stream \cite{McKayPiperno2014} contains all
unlabelled triangle-free graphs
on 15 vertices with maximum degree at most seven and 43 through 50 edges.
There are 46,184 stream records. Exact checks of $\alpha(L)\leq7$ and
the nine-set lower bound retain 997 core types, distributed as
\[
\begin{array}{c|rrrrrrrr}
m&43&44&45&46&47&48&49&50\\ \hline
\text{cores}&435&250&152&82&45&21&9&3.
\end{array}
\]
Equation \eqref{eq:d8-slack} also gives $f\leq m-36\leq14$. The shell
stream therefore contains all 6,896 unlabelled eight-vertex graphs with
7 through 14 edges. It rejects $K_4$ or an independent seven-set and,
for each remaining $F$, enumerates every weak composition of
$\sum_a\varepsilon_a\leq7$. The bound seven follows from
$f\geq7$, $m\leq50$, and \eqref{eq:d8-slack}. It then checks
\eqref{eq:d8-slack}, \eqref{eq:d8-pair-cover}, and
\eqref{eq:d8-triangle-cover}. Thus every possible fixed pair from this
branch occurs in the two filtered streams.

If the triangle condition is temporarily omitted, the other filters leave
no shell through $m=46$. At $m=47$ they leave three book graphs: two
adjacent centers, $t=3,4,5$ common leaves, and $5-t$ isolated vertices.
The forced row-size vector is
\[
 6,6,\underbrace{2,\ldots,2}_{t\text{ times}},
 \underbrace{0,\ldots,0}_{5-t\text{ times}}.
\]
A center--center--leaf triangle has row-size sum 14, contrary to
\eqref{eq:d8-triangle-cover}. Hence no case with $m\leq47$ remains.
For $m=48,49,50$, respectively, the shell filter retains 17, 40, and
48 records, while the core filter retains 21, 9, and 3 types. It follows
that the finite index contains exactly
\[
 21\cdot17+9\cdot40+3\cdot48=861
\]
core--shell pairs.

For a fixed pair, the 120 Boolean variables record membership
$b\in D_a$. The formula enforces the global edge cap; the row bounds
$|D_a|\geq d_F(a)$; the column bounds
$|\{a:b\in D_a\}|\geq\max\{0,d_L(b)-6\}$; the mixed
independent-four exclusions; the triangle-column covers; and the mixed
$K_8$ exclusions. It also enforces the local-cap families consisting of
one shell and eight core vertices, of $v$ with six shell and two core
vertices, and of $v$ with seven shell and one core vertex. Every graph in
the branch satisfies these constraints, so unsatisfiability of this
restricted formula is enough; the encoding need not include every other
nine-set family.
The standard-library semantic verifier independently reconstructs the raw
filters, the pair order, every clause family, and every auxiliary
totalizer. All 861 reconstructed formulas are refuted by checked LRAT
proofs, as detailed in Section~\ref{sec:finite-r8}. Thus $d=8$ is also
impossible, completing the proof of \eqref{eq:r8-p324}.
\end{proof}

\subsection{Human implication chain}

Assume that every nine-set in an eight-coloring of $K_{65}$ sees all eight
colors, and choose a least color graph $G$.
\begin{enumerate}[label=\textup{R8.\arabic*},leftmargin=3.5em]
\item The least-color budget is $e(G)\leq260$.
\item The shared colored recursion excludes minimum degrees at most six,
  and \eqref{eq:degree-window} gives $\delta(G)=7$.
\item For a degree-seven vertex, its 57 nonneighbors $U$ satisfy
  \[
    \alpha(G[U])\leq7,
    \qquad e(G[U])\leq232.
  \]
\item Lemma~\ref{lem:r8-p323}, \eqref{eq:r8-p324}, and the shared
  recursion give
  \[
    P_7(57)\geq235,
    \quad P_6(49)\geq206,
    \quad P_5(41)\geq177,
    \quad P_4(33)\geq149.
  \]
\item These values exceed the successive edge upper bounds
  \[
    232,\quad204,\quad176,\quad148.
  \]
  Four disjoint target copies of $K_8$ are forced in $U$.
\item Proposition~\ref{prop:block-exclusion} forbids exactly four such
  blocks when $r=8$.
\end{enumerate}

The full residual table is
\begin{center}
\begin{tabular}{c@{\quad}c@{\quad}c@{\quad}c@{\quad}c}
\toprule
$j$&order&independence cap&edge upper&no-next-$K_8$ floor\\
\midrule
0&57&7&232&235\\
1&49&6&204&206\\
2&41&5&176&177\\
3&33&4&148&149\\
\bottomrule
\end{tabular}
\end{center}

For actual induced target graphs with clique number at most seven, importing
Lemma~\ref{lem:r8-p323} and \eqref{eq:r8-p324} into
Proposition~\ref{prop:recursive-lower-bound} gives
\begin{center}
\begin{tabular}{c@{\qquad}l@{\qquad}l}
\toprule
$a$&orders $n$&corresponding floors $P_a(n)$\\
\midrule
4&31, 32, 33&120, 134, 149\\
5&40, 41, 42, 43&163, 177, 189, 203\\
6&48, 49, 50&192, 206, 218\\
7&57&235\\
\bottomrule
\end{tabular}
\end{center}
The verifier \texttt{verify\_r8\_d7\_full\_color\_bridge.py} evaluates
the recursion and records every minimizing degree branch.

\begin{proof}[Proof of Theorem~\ref{thm:r8}]
Suppose otherwise and choose a least color graph $G$. Equation
\eqref{eq:fixed-degree-reductions} gives $\delta(G)=7$. Choose a
degree-seven vertex $v$ and let $U$ be its nonneighbors. Minimum-degree
accounting gives
\[
 |U|=57,
 \qquad \alpha(G[U])\leq7,
 \qquad e(G[U])\leq260-7-\binom72=232.
\]
After deleting $j$ disjoint target copies of $K_8$, let $R_j$ be the
residual. Then
\begin{equation}\label{eq:r8-residual-data}
 |R_j|=57-8j,
 \qquad \alpha(G[R_j])\leq7-j,
 \qquad e(G[R_j])\leq232-28j.
\end{equation}
The independence claim follows by representative selection. If $R_j$
had an independent set of order $8-j$, choose one representative from
each deleted block in sequence. Every outside vertex has at most one
target neighbor in a target $K_8$, since two such neighbors would violate
the local nine-set cap. Before each choice at most seven vertices have
been selected, so at most seven vertices of the next block are forbidden.
The result is an independent eight-set in $U$; together with $v$ it
contradicts $\alpha(G)\leq8$.

Every $R_j$ is obtained only by vertex deletion, so it remains an actual
induced subgraph of the original color graph. In particular, it retains
the full-color provenance required by Lemma~\ref{lem:r8-p323} and
Proposition~\ref{prop:recursive-lower-bound}. For $j=0,1,2,3$, the four
floors in the preceding table are strictly above the respective upper
bounds $232,204,176,148$. Therefore $R_j$ contains a target $K_8$ at
each stage. Four disjoint target copies of $K_8$ are forced in $U$,
contrary to Proposition~\ref{prop:block-exclusion}.
\end{proof}

\section{Finite verification for the eight-color case}\label{sec:finite-r8}

For the degree-five endpoint, the formula searches for a labeled graph on
18 vertices that is 7-regular, triangle-free, and has independence number
at most seven, after fixing one neighborhood by symmetry. For degree eight,
write
\[
  A=N_H(v),\quad |A|=8,
  \qquad B=V(H)\setminus(A\cup\{v\}),\quad |B|=15,
\]
and put $L=\comp{H[B]}$ and $F=\comp{H[A]}$. The human reduction filters
the unlabelled core and shell streams, then indexes every pair satisfying
the remaining edge ledger.

\begin{center}
\begin{tabular}{lrrr}
\toprule
finite branch&instances&full clauses&LRATs\\
\midrule
degree five&1&50,225&1\\
degree eight&861&52,434,647&861\\
\midrule
total&862&52,484,872&862\\
\bottomrule
\end{tabular}
\end{center}
The degree-five formula has 1,233 variables and 288 totalizer merges.
Across the degree-eight formulas, the semantic verifier reconstructs
413,938 retained clauses and 836,035 totalizer merges.

The degree-eight stream reduction is
\begin{center}
\begin{tabular}{c@{\qquad}r@{\qquad}r@{\qquad}r}
\toprule
$e(L)$&filtered core types&filtered shell records&indexed pairs\\
\midrule
48&21&17&357\\
49&9&40&360\\
50&3&48&144\\
\midrule
total&33&105&861\\
\bottomrule
\end{tabular}
\end{center}
There are no indexed pairs with $e(L)\leq47$. Before the final pairing,
nauty 2.8.9 produces 46,184 raw 15-vertex core graphs and 6,896 raw
eight-vertex shell graphs. The semantic filters retain 997 core types and
905 shell records across all edge rows. The two exact generator commands
are
{\small\begin{verbatim}
geng -q -t -D7 15 43:50
geng -q 8 7:14
\end{verbatim}}
The local \texttt{geng} binary used to create the retained streams has
SHA-256
\begin{center}
\texttt{3ca950af2145c546f9f586cf960eaf98}\par\vspace{-2pt}
\texttt{f88fc3920564338f8306b6f58d018af5}.
\end{center}

The accompanying release contains the degree-five semantic verifier
\path{code/verify_endpoint18_certificate.py}, its compressed CNF and LRAT in
\path{artifacts/endpoint18/}, and the complete degree-eight package under
\path{artifacts/r8-fixed-outer/}.

The degree-eight portable archive has 2,600 regular files and a
2,599-entry checksum ledger. Its size is 15,226,513 bytes and its SHA-256
is
\begin{center}
\texttt{af5e7d7c20ab44f2330e7b69fdffe16c}\par\vspace{-2pt}
\texttt{018eb948bfb41539ff2f0a7f2e14bfc0}.
\end{center}
\section{Replay and trust boundary}\label{sec:replay}

From the \path{r7-r8/} release directory, the seven-color replay surface is
{\small\begin{verbatim}
python3 code/verify_r7_weighted_light_edge.py
sage code/verify_r7_weighted_light_edge_sage.py
python3 code/verify_r7_fixed_outer_closure.py
python3 code/verify_colored_core_ladder.py
\end{verbatim}}
The pinned LRAT checker and the full eight-color replay can be built and run
with
{\small\begin{verbatim}
git clone https://github.com/marijnheule/drat-trim.git
git -C drat-trim checkout 2e3b2dc0ecf938addbd779d42877b6ed69d9a985
make -C drat-trim lrat-check
./reproduce.sh --lrat-check "$PWD/drat-trim/lrat-check"
\end{verbatim}}

The LRAT format and its checking semantics are described in
\cite{CruzFilipeEtAl2017}; the DRAT-trim toolchain is described in
\cite{WetzlerHeuleHunt2014}. The audited \texttt{lrat-check} executable was
built from DRAT-trim source commit
\begin{center}
\texttt{2e3b2dc0ecf938addbd779d42877b6ed69d9a985}.
\end{center}
Its arm64 binary SHA-256 is
\begin{center}
\texttt{a2caddba197bbb0846aabf2dd54e87d7}\par\vspace{-2pt}
\texttt{9612e3b840ea0c43714f5c522f2d86c7}.
\end{center}

The finite checkers do not verify the human implication chains, the
Kang--Pikhurko theorem, or completeness of nauty's unlabelled graph
generation. For $r=8$, the independent semantic checkers establish what
the retained formulas mean; the LRAT checker establishes their
unsatisfiability. The human reduction is needed to connect those formulas
to Theorem~\ref{thm:r8}. For $r=7$, the proof receipt is direct exhaustive
enumeration plus an unlabelled reconstruction, not a detached certificate.

Certificate payloads are in the artifact directories named in
Sections~\ref{sec:finite-r7} and~\ref{sec:finite-r8}.

\section{AI-use disclosure}\label{sec:ai-disclosure}

OpenAI GPT 5.6 Sol was used substantially in this work. Its use included
exploratory proof search, proposing intermediate lemmas, checking algebra,
drafting and revising verification programs, organizing finite case
reductions, preparing replay ledgers, and drafting portions of the research
notes from which this manuscript was assembled. The model
also helped identify failed approaches and presentation errors. This was
not limited to copy editing.

OpenAI GPT-5 (Codex) was used for later source review, mathematical audits,
certificate replay, and preparation of the release package. xAI Grok 4.5
was used to critique a draft.

Model output is not treated as mathematical authority. Proposed lemmas were
rejected when their hypotheses could not be traced, and finite claims were
kept separate from human reductions. The $r=7$ finite lemma has two
implementations with different graph-generation methods. The $r=8$
formulas have separate production and semantic-verification paths, followed
by LRAT replay. These checks reduce some risks, but they do not amount to an
independent human review and do not certify the cited published theorem.

This preprint is being circulated for independent checking and has not yet
received external mathematical review. The author remains responsible for
the statements, proof, code, citations, and any errors.

Any circulated revision should retain a clear disclosure of the model's
substantial role and link to the human-readable proof, semantic verifiers,
manifests, and certificate replay instructions. The disclosure should not
be read as a claim that the model independently proved or formally verified
the results.

\section{Nonclaims and remaining work}\label{sec:nonclaims}

This preprint makes none of the following claims.
\begin{enumerate}[label=\textup{N\arabic*.},leftmargin=2.8em]
\item It does not solve Erd\H{o}s Problem 617 for arbitrary $r$.
\item It does not settle the fixed case $r=9$ or any larger unresolved
  fixed parameter.
\item It does not prove the cited Kang--Pikhurko theorem from first
  principles.
\item It is not a Lean formalization.
\item It has not been externally reviewed or accepted.
\item A passing finite replay does not replace the human reductions that
  establish completeness of the case splits.
\end{enumerate}

The older degree-nine $\mathsf Q_3$, $\mathsf P(8,4)$, and
$\mathsf P(8,5)$ statements remain open as standalone claims. The fixed
$r=8$ bridge bypasses them rather than proving them.

\appendix
\section{Dependency ledger}

\begin{center}
\begin{tabular}{L{0.25\textwidth}L{0.31\textwidth}L{0.32\textwidth}}
\toprule
result&human dependency&finite dependency\\
\midrule
Theorem~\ref{thm:r7}&shared colored ladder; star-and-cover propagation;
three-block contradiction&weighted light-edge enumeration and unlabelled
reconstruction\\
Theorem~\ref{thm:r8}&shared colored ladder; 23-vertex equality-core proof;
human branches of $P_3(24)$; four-block contradiction&one endpoint formula
and 861 core--shell formulas, semantic reconstruction, 862 LRAT replays\\
\bottomrule
\end{tabular}
\end{center}

The external mathematical inputs are Brooks's theorem, the
Andr\'asfai--Erd\H{o}s--S\'os minimum-degree theorem, and the
Kang--Pikhurko extremal theorem, including its equality description in
the instances explicitly cited above. The accompanying README identifies
the source, code, proof artifacts, and replay procedure.

\bibliographystyle{amsplain}
\bibliography{references}

\end{document}
